\documentclass[a4paper,12pt]{article}

\usepackage[ngerman]{babel}
\usepackage[utf8]{inputenc}
\usepackage[T1]{fontenc}
\usepackage{lmodern}
\usepackage{geometry}
\usepackage{listings}
\usepackage{xcolor}

\geometry{margin=2.5cm}

\title{Praktische Übung: Systemkomponenten für eine Webanwendung auswählen}
\author{Modul 321}
\date{\today}

\lstset{
    basicstyle=\ttfamily\small,
    breaklines=true,
    frame=single,
    backgroundcolor=\color{gray!10}
}

\begin{document}

\maketitle

\section*{Ziel der Übung}

In dieser Übung plant ihr für eine einfache Webanwendung passende
\textbf{Systemkomponenten}. Ziel ist nicht eine vollständige Produktentwicklung,
sondern das saubere Einordnen technischer Verantwortungen.

\section*{Ausgangslage}

Ihr entwickelt eine kleine Lernplattform mit diesen Anforderungen:

\begin{itemize}
\item Benutzer melden sich an
\item Kurse werden angezeigt
\item Dateien können hochgeladen werden
\item Nach einem Upload sollen andere Teile informiert werden
\item Das System soll bei vielen Anfragen stabil bleiben
\item Fehler und Ausfälle sollen sichtbar sein
\end{itemize}

\section*{Aufgabe}

Entwerft eine einfache Architektur mit passenden Systemkomponenten.

Mindestens diese Kategorien sollen vorkommen:

\begin{itemize}
\item \textbf{Identitätsverwaltung}
\item \textbf{Ereignisverwaltung}
\item \textbf{Monitoring}
\item \textbf{Datenhaltung}
\item \textbf{Lastenverteilung}
\end{itemize}

\section*{Variante A: Skizze ohne Code}

Erstellt ein Diagramm oder ein Dokument mit:

\begin{itemize}
\item Frontend
\item Backend / API
\item den gewählten Systemkomponenten
\item kurzen Pfeilen oder Beschriftungen, wie die Teile zusammenarbeiten
\end{itemize}

\section*{Variante B: Kleine technische Demo}

Erstellt eine einfache Ordner- oder Codestruktur, zum Beispiel:

\begin{lstlisting}
system-components-demo/
|
|-- auth/
|   `-- login.js
|
|-- events/
|   `-- event-bus.js
|
|-- monitoring/
|   `-- logger.js
|
|-- data/
|   `-- repository.js
|
|-- loadbalancer/
|   `-- router.js
|
`-- app.js
\end{lstlisting}

Die Dateien dürfen sehr klein bleiben. Es reicht, wenn klar wird, welche
Verantwortung jede Komponente hat.

\section*{Arbeitsaufträge}

\subsection*{Aufgabe 1: Komponenten benennen}
Nennen Sie zu jeder der fünf Kategorien:
\begin{enumerate}
\item die Aufgabe der Komponente
\item ein Beispiel aus der Lernplattform
\item warum diese Aufgabe nicht einfach in einer beliebigen anderen Komponente liegen sollte
\end{enumerate}

\subsection*{Aufgabe 2: Systemfluss beschreiben}
Beschreiben Sie kurz einen möglichen Ablauf:

\begin{enumerate}
\item Benutzer meldet sich an
\item Datei wird hochgeladen
\item ein Ereignis wird ausgelöst
\item das System protokolliert den Vorgang
\end{enumerate}

\subsection*{Aufgabe 3: Last und Betrieb}
Beschreiben Sie:
\begin{itemize}
\item wie Lastenverteilung helfen würde
\item wie Monitoring helfen würde
\item welche Probleme ohne diese Komponenten entstehen könnten
\end{itemize}

\subsection*{Aufgabe 4: Risiken bei falscher Zuordnung}
Beantworten Sie:
\begin{itemize}
\item Was wäre problematisch, wenn Login-Logik direkt überall im Code verteilt wäre?
\item Was wäre problematisch, wenn es keine Ereignisverwaltung gäbe?
\item Was wäre problematisch, wenn Monitoring fehlt?
\end{itemize}

\subsection*{Aufgabe 5: Minimalarchitektur}
Zeichnen oder beschreiben Sie eine einfache Zielarchitektur in 5 bis 8 Schritten.

\section*{Abgabe}

Die Abgabe soll enthalten:

\begin{itemize}
\item kurze Beschreibung der Anwendung
\item Übersicht über die gewählten Komponenten
\item Skizze, Diagramm oder kleine Code-Demo
\item kurze Begründung der Rollen der Komponenten
\item kurze Reflexion über Risiken oder fehlende Teile
\end{itemize}

\section*{Reflexionsfragen}

\begin{enumerate}
\item Welche Komponente wäre aus Ihrer Sicht am sicherheitskritischsten?
\item Welche Komponente wäre für Betrieb und Stabilität besonders wichtig?
\item Welche Komponente hilft, das System später leichter zu erweitern?
\item Wo wäre zu viel Komplexität für ein kleines Projekt eher unnötig?
\end{enumerate}

\section*{Bewertung}

\begin{itemize}
\item Kategorien sind korrekt erkannt
\item Aufgaben der Komponenten sind nachvollziehbar beschrieben
\item Zusammenspiel der Komponenten ist verständlich
\item Risiken oder Fehlannahmen werden erkannt
\item Die Lösung bleibt einfach und trotzdem fachlich sauber
\end{itemize}

\end{document}
